% ASSERT_CONVENTION: natural_units=dimensionless, metric_signature=mostly_minus, gamma_matrix_convention=Cl(9,0), generator_normalization=gamma_ab/4, commutation_convention=[A,B]=AB-BA
%
% Phase 48: Lorentz Equivariance Derivation Summary
% Produced by Plan 48-02
% Date: 2026-04-12

\documentclass[11pt]{article}
\usepackage{amsmath,amssymb,amsthm,booktabs,geometry}
\geometry{margin=1in}

\newtheorem{theorem}{Theorem}
\newtheorem{proposition}{Proposition}
\newtheorem{corollary}{Corollary}
\theoremstyle{definition}
\newtheorem{definition}{Definition}
\theoremstyle{remark}
\newtheorem{remark}{Remark}

\title{Stabilizer Structure and Lorentz Equivariance\\ of the $\pi_u$ Projection on $V_0$}
\author{Phase 48}
\date{April 2026}

\begin{document}
\maketitle

%% ============================================================
\section{Setup}
%% ============================================================

Let $V_0 = h_2(\mathbb{O}) \cong \mathbb{R}^{10}$ be the Peirce-0 eigenspace
of $h_3(\mathbb{O})$ with respect to the idempotent $E_{11}$.  The group
$\mathrm{Spin}(9)$ acts on $V_0$ via the 10-dimensional (vector)
representation derived from the commutator action
\begin{equation}\label{eq:comm-action}
  \bigl[\tfrac{1}{4}\gamma_a\gamma_b,\; T_c\bigr]
  = \sum_{d=0}^{9} M^{(ab)}_{dc}\, T_d,
  \qquad 0 \le a < b \le 8,
\end{equation}
where $\gamma_a$ ($a=0,\ldots,8$) are the $\mathrm{Cl}(9,0)$ generators
satisfying $\{\gamma_a,\gamma_b\} = 2\delta_{ab}\,I_{16}$,
and $\{T_c\}_{c=0}^{9}$ is the $V_0$ basis from Phase~46.

Choosing the unit imaginary octonion $u = e_7$ defines the projection
$\pi_u : V_0 \to h_2(\mathbb{C}_u)$, whose image is the 4-dimensional
spacetime factor $h_2(\mathbb{C}_u)$ with Lorentzian metric
$\det_2 \leftrightarrow \eta = \mathrm{diag}(+1,-1,-1,-1)$
(Phase~46, Eq.~46.2).

The $V_0$ coordinate basis splits as:
\begin{itemize}
  \item \textbf{Spacetime indices} $S = \{0,1,2,9\}$:
    $c_0 = \beta+\gamma$,\; $c_1 = \beta-\gamma$,\;
    $c_2 = \mathrm{Re}(x_1)$,\; $c_9 = (x_1)_{e_7}$.
  \item \textbf{Internal indices} $I = \{3,4,5,6,7,8\}$:
    the six $W$-sector components of $x_1$, where
    $W = u^\perp \cap \mathrm{Im}(\mathbb{O}) = \mathrm{span}\{e_1,\ldots,e_6\}$.
\end{itemize}

%% ============================================================
\section{$V_0$ Stabilizer: $\mathfrak{so}(3) \oplus \mathfrak{so}(6)$}
%% ============================================================

\begin{definition}
The \emph{$V_0$ stabilizer} is
\[
  \mathfrak{stab}(u)
  = \bigl\{L \in \mathfrak{spin}(9) :
    L \text{ preserves the } 4{+}6 \text{ splitting}\bigr\},
\]
equivalently, the subalgebra whose 10$\times$10 matrices are block-diagonal
with respect to $(S, I)$.
\end{definition}

\begin{theorem}[Stabilizer Structure, Phase 48-01]\label{thm:stab}
\[
  \mathfrak{stab}(u) \;=\; \mathfrak{so}(3) \oplus \mathfrak{so}(6),
  \qquad \dim = 3 + 15 = 18.
\]
The $\mathfrak{so}(3)$ factor acts nontrivially on the 4-dim spacetime
block and trivially on the 6-dim internal block.
The $\mathfrak{so}(6)$ factor acts nontrivially on the internal block
and trivially on spacetime.  Cross-brackets vanish:
$[\mathfrak{so}(3), \mathfrak{so}(6)] = 0$.
\end{theorem}

\noindent\textbf{Verification} (computational, exact to machine precision):
\begin{itemize}
  \item SVD of the 48$\times$36 off-diagonal constraint matrix:
    18 null singular values (all $< 10^{-17}$), gap 0.71.
  \item Stabilizer closure: 153 brackets, max residual $5.3 \times 10^{-17}$.
  \item Killing form: $\mathrm{spec}(K) = \{-2^{(\times 15)},\, -0.5^{(\times 3)}\}$,
    negative definite, center dimension 0.
\end{itemize}

\begin{remark}
The stabilizer dimension is 18, not 21.
One might expect $\dim\,\mathfrak{so}(3,1) + \dim\,\mathfrak{so}(6) = 6 + 15 = 21$,
but $\mathfrak{so}(3,1)$ is \emph{non-compact} and cannot embed as a
subalgebra of the compact $\mathfrak{so}(9)$.
Only its maximal compact subalgebra $\mathfrak{so}(3)$ (spatial rotations)
embeds.  The 3 boost generators do not exist within $\mathfrak{spin}(9)$.
\end{remark}

%% ============================================================
\section{Lorentz Rotation Subalgebra}
%% ============================================================

The Minkowski basis for $h_2(\mathbb{C}_u)$ is
$\{x_0, x_1, x_2, x_3\}$ with $x_0 = c_0/2$, $x_1 = c_2$,
$x_2 = c_9$, $x_3 = c_1/2$.  The basis transformation is
\begin{equation}\label{eq:basis-change}
  B = \begin{pmatrix} 1/2 & 0 & 0 & 0 \\
                       0 & 0 & 1 & 0 \\
                       0 & 0 & 0 & 1 \\
                       0 & 1/2 & 0 & 0 \end{pmatrix},
  \qquad
  L_{\mathrm{Mink}} = B\, L_S\, B^{-1}.
\end{equation}

In the Minkowski basis, the 3 rotation generators are:
\begin{equation}\label{eq:rotations}
  J_{13} = \begin{pmatrix} 0&0&0&0\\0&0&0&-\tfrac12\\0&0&0&0\\0&\tfrac12&0&0 \end{pmatrix},
  \quad
  J_{23} = \begin{pmatrix} 0&0&0&0\\0&0&0&0\\0&0&0&-\tfrac12\\0&0&\tfrac12&0 \end{pmatrix},
  \quad
  J_{12} = \begin{pmatrix} 0&0&0&0\\0&0&\tfrac12&0\\0&-\tfrac12&0&0\\0&0&0&0 \end{pmatrix}.
\end{equation}
These satisfy the $\mathfrak{so}(3)$ commutation relations
\begin{equation}\label{eq:so3-comm}
  [J_{13}, J_{23}] = -\tfrac{1}{2} J_{12}, \quad
  [J_{13}, J_{12}] = +\tfrac{1}{2} J_{23}, \quad
  [J_{23}, J_{12}] = -\tfrac{1}{2} J_{13},
\end{equation}
i.e.\ $[J_i, J_j] = \tfrac{1}{2}\epsilon_{ijk}\, J_k$ with the
generator normalization from $\gamma_{ab}/4$.

\begin{proposition}[Metric compatibility]\label{prop:eta}
Every rotation generator satisfies
\begin{equation}\label{eq:eta-compat}
  \eta\, L_{\mathrm{Mink}} + L_{\mathrm{Mink}}^{\!\top}\, \eta = 0,
  \qquad \eta = \mathrm{diag}(+1,-1,-1,-1),
\end{equation}
with max error $2.2 \times 10^{-16}$.
\end{proposition}

\noindent This is the defining relation of $\mathfrak{so}(3,1)$
generators.  The 3 rotation generators satisfy it because
spatial rotations preserve both the Euclidean and Lorentzian metrics.
They form the \emph{maximal compact subalgebra} of $\mathfrak{so}(3,1)$.

\begin{remark}[Compact vs.\ non-compact]
The abstract Lie algebra of the spacetime block is $\mathfrak{so}(3)$,
which is simultaneously:
\begin{itemize}
  \item a subalgebra of $\mathfrak{so}(4)$ (compact, Euclidean rotations),
  \item the rotation subalgebra of $\mathfrak{so}(3,1)$ (Lorentz algebra),
  \item the maximal compact subalgebra of $\mathfrak{so}(3,1)$.
\end{itemize}
The Lorentz boosts ($K_1, K_2, K_3$) would need non-compact generators
that do not exist in the compact $\mathfrak{spin}(9)$.
At the group level: $\mathrm{SU}(2) \subset \mathrm{Spin}(9)$ acts as
spatial rotations $\mathrm{SO}(3) \subset \mathrm{SO}(3,1)$ on
$(h_2(\mathbb{C}_u), \det_2)$.  The full Lorentz group
$\mathrm{SL}(2,\mathbb{C})$ does not embed in compact $\mathrm{Spin}(9)$.
This is the standard Baez~(2002) isomorphism
$\mathrm{SL}(2,\mathbb{K}) = \mathrm{Spin}(\dim\mathbb{K}+1, 1)$
restricted to the compact subgroup.
\end{remark}

%% ============================================================
\section{Internal Symmetry: $\mathfrak{so}(6) \cong \mathfrak{su}(4)$}
%% ============================================================

The 15 internal generators act on the 6-dim $W$-sector with:
\begin{itemize}
  \item Killing form: all eigenvalues $= -2$ (15-fold degenerate),
    negative definite $\Rightarrow$ compact real form.
  \item $\mathfrak{so}(6) \cong \mathfrak{su}(4)$
    (the standard isomorphism of rank-3 algebras).
\end{itemize}

The relationship to the Standard Model gauge algebra $G_{\mathrm{SM}}$
(dim~8, from the $V_{1/2}$ stabilizer of $J_u$, Krasnov 2019) is:
$G_{\mathrm{SM}} \subset \mathfrak{stab}(u)$, verified with
max residual $4.85 \times 10^{-15}$.
The extra $15 - 8 = 7$ generators in $\mathfrak{so}(6)$ beyond
$G_{\mathrm{SM}}$ preserve the $V_0$ splitting but do NOT commute
with $J_u$ on $V_{1/2}$.

%% ============================================================
\section{Equivariance Theorem}
%% ============================================================

\begin{theorem}[Equivariance of $\pi_u$]\label{thm:equiv}
For all $L \in \mathfrak{stab}(u)$ and all $Y \in V_0$:
\begin{equation}\label{eq:equivariance}
  \pi_u(L \cdot Y) = L\big|_{h_2(\mathbb{C}_u)} \cdot \pi_u(Y).
\end{equation}
\end{theorem}

\begin{proof}
Let $P_S : V_0 \to h_2(\mathbb{C}_u)$ be the projection onto spacetime
indices $S = \{0,1,2,9\}$.  Since $L \in \mathfrak{stab}(u)$, the
10$\times$10 matrix of $L$ is block-diagonal with respect to $(S, I)$:
$L = L_S \oplus L_I$.  Therefore
\[
  P_S \circ L = L_S \circ P_S,
\]
which is precisely the equivariance condition~\eqref{eq:equivariance}
since $\pi_u = P_S$ on $V_0$.

Verified numerically: $\max_{L,\, e_k}
|P_S(L \cdot e_k) - L_S(P_S \cdot e_k)| = 5.2 \times 10^{-17}$
over all 18 stabilizer generators and all 10 basis vectors.
\end{proof}

%% ============================================================
\section{Mixing Generators and Physical Interpretation}
%% ============================================================

The 18 coset generators (mixing generators) in $\mathfrak{spin}(9) / \mathfrak{stab}(u)$
have nonzero off-diagonal blocks: they mix spacetime and internal
components.  These correspond to infinitesimal changes of $u$ in $S^6$
(equivalently, changes of complex structure on $\mathbb{O}$).

\begin{center}
\begin{tabular}{lccp{6cm}}
\toprule
\textbf{Subalgebra} & \textbf{Dim} & \textbf{Type} & \textbf{Physical Role} \\
\midrule
$\mathfrak{stab}(u)$ & 18 & $\mathfrak{so}(3) \oplus \mathfrak{so}(6)$
  & Symmetry preserving the $u$-splitting \\
\quad Rotations & 3 & $\mathfrak{so}(3)$
  & Spatial rotations on $h_2(\mathbb{C}_u)$; maximal compact
    subalgebra of $\mathfrak{so}(3,1)$ \\
\quad Internal & 15 & $\mathfrak{so}(6) \cong \mathfrak{su}(4)$
  & Internal symmetry of $W$-sector \\
$G_{\mathrm{SM}}$ & 8 & $\mathfrak{su}(3) \oplus \mathfrak{u}(1) \oplus \cdots$
  & Standard Model gauge algebra (from $V_{1/2}$; $\subset\!\mathfrak{stab}(u)$) \\
Mixing & 18 & coset
  & Transformations changing $u \in S^6$ \\
\bottomrule
\end{tabular}
\end{center}

\bigskip
\noindent\textbf{Dimension accounting:}
\[
  \underbrace{3}_{\text{rotations}} + \underbrace{15}_{\text{internal}}
  = \underbrace{18}_{\text{stabilizer}},
  \qquad
  \underbrace{18}_{\text{stabilizer}} + \underbrace{18}_{\text{mixing}}
  = \underbrace{36}_{\mathfrak{spin}(9)}.
\]

%% ============================================================
\section*{References}
%% ============================================================

\begin{itemize}
  \item Baez, J.C.\ (2002).
    The octonions. \emph{Bull.\ AMS} \textbf{39}, 145--205.
    arXiv:math/0105155, Section~3.3.
  \item Krasnov, K.\ (2019).
    arXiv:1912.11282, Section~3.
  \item Phase~46: $\pi_u$ projection, $\det_2$ Gram matrix, spacetime/internal splitting.
  \item Phase~47: $\mathrm{Spin}(9)$ action on $V_0$ verified (630 tests), uniqueness.
  \item Phase~48-01: $V_0$ stabilizer = $\mathfrak{so}(3) \oplus \mathfrak{so}(6)$, dim 18.
\end{itemize}

\end{document}
